\section*{Capítulo \ref{ch:b1:periodicity} --- Periodicidade}

\begin{solution}{exo:b1:periodicity:1}
\ce{Ca} $[\ce{Ar}]\,4\text{s}^2$: \omterm{def:b1:periodicity:table}{período} 4, \omterm{def:b1:periodicity:table}{grupo} 2, \omterm{def:b1:periodicity:table}{bloco} s. \ce{Fe}
$[\ce{Ar}]\,3\text{d}^6 4\text{s}^2$: \omterm{def:b1:periodicity:table}{período} 4, \omterm{def:b1:periodicity:table}{grupo} 8, \omterm{def:b1:periodicity:table}{bloco} d. \ce{Se}
$[\ce{Ar}]\,3\text{d}^{10} 4\text{s}^2 4\text{p}^4$: \omterm{def:b1:periodicity:table}{período} 4, \omterm{def:b1:periodicity:table}{grupo} 16, \omterm{def:b1:periodicity:table}{bloco}
p. \ce{I} $[\ce{Kr}]\,4\text{d}^{10} 5\text{s}^2 5\text{p}^5$: \omterm{def:b1:periodicity:table}{período} 5,
\omterm{def:b1:periodicity:table}{grupo} 17, \omterm{def:b1:periodicity:table}{bloco} p.
\end{solution}

\begin{solution}{exo:b1:periodicity:2}
\ce{Na} é maior que \ce{Cl}: mesmo \omterm{def:b1:periodicity:table}{período}, $Z^{*}$ menor. \ce{Na} é
maior que \ce{Na+}, que perdeu sua camada 3s. \ce{Cl-}
(\qty{181}{pm}) é maior que \ce{F-} (\qty{133}{pm}): uma camada a mais.
\ce{Na+} (\qty{102}{pm}) é maior que \ce{Mg^{2+}} (\qty{72}{pm}): os mesmos
dez elétrons, carga nuclear 11 contra 12.
% ledger: rion:Cl-VI, rion:F-VI, rion:Na+VI, rion:Mg2+VI
\end{solution}

\begin{solution}{exo:b1:periodicity:3}
$\ce{Na} (5.14) < \ce{Al} (5.99) < \ce{Mg} (7.65) < \ce{S} (10.36) <
\ce{P} (10.49) < \ce{Ar} (15.76)$. O aumento geral ao longo do \omterm{def:b1:periodicity:table}{período}
é quebrado duas vezes: \ce{Al} abaixo de \ce{Mg} (o elétron retirado é o
primeiro elétron 3p, mais alto em energia que 3s) e \ce{S} abaixo de \ce{P}
(o elétron retirado é o primeiro elétron 3p emparelhado).
\end{solution}

\begin{solution}{exo:b1:periodicity:4}
A \omterm{def:b1:periodicity:polarisability}{polarizabilidade} mede a facilidade com que a nuvem eletrônica é
deformada por um campo elétrico. \ce{I-} é mais polarizável que \ce{F-}:
é maior e seus elétrons externos estão mais longe do núcleo. \ce{Cl-} (18
elétrons retidos por $Z = 17$) é muito mais polarizável que \ce{Na+} (10
elétrons retidos por $Z = 11$), como os ânions em relação aos cátions.
\end{solution}

\begin{solution}{exo:b1:periodicity:5}
Razões: $18.83/5.99 = 3.1$, $28.45/18.83 = 1.5$, $119.99/28.45 = 4.2$. O
salto vem depois do terceiro elétron: três \omterm{def:b1:quantum-numbers:configuration}{elétrons de valência}, \omterm{def:b1:periodicity:table}{grupo}
13, \omterm{def:b1:periodicity:table}{período} 3: alumínio. Ele forma \ce{Al^{3+}}, $1\text{s}^2 2\text{s}^2
2\text{p}^6$.
% ledger: ie:Al, ie2:Al, ie3:Al, ie4:Al
\end{solution}

\begin{solution}{exo:b1:periodicity:6}
O elétron acrescentado em \ce{F-} entra na pequena camada 2p, onde é
fortemente repelido pelos sete \omterm{def:b1:quantum-numbers:configuration}{elétrons de valência} que já estão lá; na
camada 3p, maior, do cloro, a repulsão é menor, de modo que se libera mais
energia. No nitrogênio ($2\text{p}^3$, \omorbs{u,u,u}), o elétron extra
teria de se emparelhar em uma \omterm{def:b1:quantum-numbers:orbital}{subcamada} semipreenchida; no magnésio
($3\text{s}^2$), teria de iniciar a \omterm{def:b1:quantum-numbers:orbital}{subcamada} 3p, mais alta. Em nenhum dos
dois casos o ânion é mais estável que o átomo e um elétron livre.
\end{solution}

\begin{solution}{exo:b1:periodicity:7}
$\chi_M(\ce{Na}) = (5.14 + 0.55)/2 = \qty{2.85}{eV}$; $\chi_M(\ce{Cl}) =
(12.97 + 3.61)/2 = \qty{8.29}{eV}$. A diferença é grande: o par de
elétrons de uma ligação \ce{Na-Cl} pertence quase inteiramente ao cloro,
e o cloreto de sódio é iônico, \ce{Na+ Cl-}.
\end{solution}

\begin{solution}{exo:b1:periodicity:8}
\ce{H-F}: $\Delta\chi = 1.78$, no limite da região iônica (na prática,
uma ligação covalente muito polar), \ce{F^{\delta-}}. \ce{C-H}: 0.35,
quase apolar. \ce{Na-Cl}: 2.23, iônica, \ce{Cl^{-}}. \ce{C-Cl}: 0.61,
covalente polar, \ce{C^{\delta+}-Cl^{\delta-}}. \ce{O-H}: 1.24, covalente
polar, \ce{O^{\delta-}-H^{\delta+}}.
\end{solution}

\begin{solution}{exo:b1:periodicity:9}
Descendo o grupo, o \omterm{def:b1:quantum-numbers:configuration}{elétron de valência} está em uma camada de $n$ maior,
mais longe do núcleo e blindado por mais \omterm{def:b1:quantum-numbers:configuration}{elétrons de caroço}; ele é
retirado com mais facilidade. O rubídio, uma camada além, tem \omterm{def:b1:periodicity:ionisation-energy}{primeira
energia de ionização} abaixo de
\qty{4.34}{eV}.
\end{solution}

\begin{solution}{exo:b1:periodicity:10}
O zinco é $[\ce{Ar}]\,3\text{d}^{10} 4\text{s}^2$: seu elétron vem de uma
\omterm{def:b1:quantum-numbers:orbital}{subcamada} 4s completa; o gálio, $\dots 4\text{s}^2 4\text{p}^1$, perde seu
único elétron 4p, mais alto em energia e blindado pelas \omterm{def:b1:quantum-numbers:orbital}{subcamadas} 4s e 3d
completas. O arsênio, $4\text{p}^3$, perde um elétron de uma \omterm{def:b1:quantum-numbers:orbital}{subcamada}
semipreenchida; o selênio, $4\text{p}^4$, perde o elétron emparelhado,
empurrado para cima pela repulsão de seu parceiro: sua energia de
ionização é ligeiramente menor.
\end{solution}

\begin{solution}{exo:b1:periodicity:11}
$\Delta E = (3.98 - 2.20)^2 = 1.78^2 = \qty{3.17}{eV} = 3.17 \times 96.5 =
\qty{306}{kJ/mol}$. Então $D(\ce{H-F}) = \tfrac12(436 + 158) + 306 = 297 +
306 \approx \qty{603}{kJ/mol}$. O grande excesso é a atração adicional
entre as cargas parciais $\ce{H^{\delta+}}$ e $\ce{F^{\delta-}}$, devida
à grande diferença de \omterm{def:b1:periodicity:electronegativity}{eletronegatividade}.
\end{solution}

\begin{solution}{exo:b1:periodicity:12}
A \omterm{def:b1:periodicity:electronegativity}{escala de Pauling} é definida a partir das energias das ligações A--B:
um elemento que não forma nenhuma ligação não tem valor de Pauling. Hélio,
neônio e argônio não formam compostos estáveis. O criptônio e sobretudo o
xenônio têm camadas de valência maiores e mais polarizáveis e energias de
ionização menores (\qty{14.0}{eV} para o criptônio contra \qty{15.76}{eV}
para o argônio); eles são oxidados pelo flúor e pelo oxigênio e formam
compostos como \ce{XeF2}, de modo que suas energias de ligação podem ser
medidas.
% ledger: ie:Kr, ie:Ar
\end{solution}

\begin{solution}{pb:b1:periodicity:1}
\textbf{1.} Cada elétron é retirado de um íon mais positivo e menor que o
anterior, que retém seus elétrons com mais força.
\textbf{2.} $15.04/7.65 = 1.97$; $80.14/15.04 = 5.33$; $109.27/80.14 = 1.36$.
\textbf{3.} O salto vem depois de dois elétrons: X tem dois \omterm{def:b1:quantum-numbers:configuration}{elétrons de
valência}, \omterm{def:b1:periodicity:table}{grupo} 2; no \omterm{def:b1:periodicity:table}{período} 3, é o magnésio.
\textbf{4.} \ce{Mg^{2+}}, $1\text{s}^2 2\text{s}^2 2\text{p}^6$, a
configuração do neônio.
\textbf{5.} Um metal: \omterm{def:b1:periodicity:table}{grupo} 2, à esquerda da tabela, com baixas \omterm{def:b1:periodicity:ionisation-energy}{primeiras
energias de ionização}.
\textbf{6.} O terceiro elétron vem do cerne 2p, de $n$ menor, muito mais
próximo do núcleo e quase sem \omterm{def:b1:periodicity:screening}{blindagem}, e é retirado de um íon com
carga dupla.
% ledger: ie:Mg, ie2:Mg, ie3:Mg, ie4:Mg
\textbf{7.} Todos têm dez elétrons: $1\text{s}^2 2\text{s}^2 2\text{p}^6$.
\textbf{8.} Os mesmos dez elétrons são retidos por cargas nucleares 8, 9,
11, 12, 13: a nuvem se contrai à medida que $Z$ cresce.
\textbf{9.} \ce{O^{2-}}: o maior, com a menor carga nuclear para seus
dez elétrons, e um ânion.
\textbf{10.} $140/54 = 2.6$.
\textbf{11.} $198.8/2 = \qty{99.4}{pm}$; o ânion é $181/99.4 = 1.8$
vezes maior que o \omterm{def:b1:periodicity:radius}{raio covalente}.
% ledger: re:Cl2, rion:Cl-VI
\textbf{12.} $\ce{Mg^{2+}} < \ce{Na+} < \ce{Cl-}$: os dois cátions são
isoeletrônicos (carga 12 contra 11), e \ce{Cl-} tem uma terceira camada.
\textbf{13.} $\chi_M$ = 10.41 (\ce{F}), 8.29 (\ce{Cl}), 7.59 (\ce{Br}),
7.53 (\ce{O}), 6.26 (\ce{C}), em eV.
\textbf{14.} Mulliken: $\ce{F} > \ce{Cl} > \ce{Br} > \ce{O} > \ce{C}$;
Pauling: $\ce{F} > \ce{O} > \ce{Cl} > \ce{Br} > \ce{C}$. O oxigênio passa
da quarta para a segunda posição.
\textbf{15.} $3.98/10.41 = 0.382$, $3.16/8.29 = 0.381$, $2.96/7.59 =
0.390$: quase constante, cerca de 0.38; para os halogênios, as duas
escalas são proporcionais.
\textbf{16.} A \omterm{def:b1:periodicity:electron-affinity}{afinidade eletrônica} do oxigênio é pequena (\qty{1.44}{eV})
porque o elétron acrescentado precisa se emparelhar na compacta camada
2p; o valor de Mulliken do átomo livre subestima a atração que o oxigênio
exerce em suas ligações, que a \omterm{def:b1:periodicity:electronegativity}{escala de Pauling}, construída a partir das
ligações, mede diretamente.
\textbf{17.} O cloro em \ce{HCl}; o flúor em \ce{ClF}.
\textbf{18.} $\Delta\chi = 0.96$: covalente polar.
\textbf{19.} $D(\ce{H-H}) = 2 \times 218.0 = \qty{436.0}{kJ/mol}$.
\textbf{20.} $D(\ce{Cl-Cl}) = 2 \times 121.3 = \qty{242.6}{kJ/mol}$.
\textbf{21.} $D(\ce{H-Cl}) = 218.0 + 121.3 - (-92.3) = \qty{431.6}{kJ/mol}$.
\textbf{22.} $\Delta E = 431.6 - \tfrac12(436.0 + 242.6) = 431.6 - 339.3 =
\qty{92.3}{kJ/mol}$. Em símbolos, $D(\ce{H-Cl}) = \Delta_fH(\ce{H}) +
\Delta_fH(\ce{Cl}) - \Delta_fH(\ce{HCl})$, enquanto a média das ligações
simétricas é $\Delta_fH(\ce{H}) + \Delta_fH(\ce{Cl})$, de modo que $\Delta E =
-\Delta_fH(\ce{HCl})$.
\textbf{23.} $92.3/96.485 = \qty{0.957}{eV}$.
\textbf{24.} $|\chi(\ce{Cl}) - \chi(\ce{H})| = \sqrt{0.957} =
\textbf{0.98}$, contra $3.16 - 2.20 = 0.96$ nas tabelas: a aritmética de
Pauling, feita com dados modernos, reproduz a diferença tabelada com
precisão de 0.02.
\end{solution}
